\begin{titlepage}
\thispagestyle{empty}
\vspace*{1.0in}
{\color{accent}\large RESEARCH MANUSCRIPT\par}
\vspace{0.4in}
{\color{ink}\fontsize{34}{39}\selectfont\bfseries Asset Price Bubbles\par}
\vspace{0.22in}
{\color{ink}\LARGE Theory, Signals, and Traded Contracts\par}
\vspace{0.5in}
{\large Rebalancing and identification\par
Daily and intraday monitoring\par
Leveraged funds, options, and constituents\par}
\vfill
\noindent A unified account of the research, with economic explanations,
principal mathematical results, and evidence from simulations and market data.
\par\bigskip
\noindent Draft prepared 9 October 2026\par
\noindent Evidence through the research snapshot documented in Appendix~\ref{app:record}.
\end{titlepage}
\pagenumbering{roman}
\markboth{Overview}{}
\section*{Abstract}
\phantomsection
\addcontentsline{toc}{section}{Abstract}
In the strict-local-martingale theory of asset prices, a bubble is a
shortfall of expected future value under a pricing measure relative to the
current price. This manuscript studies how that valuation concept relates
to patterns in observed prices and to the wealth produced by trading
strategies. We classify constant-proportion wealth by exposure and the
asset's behavior near zero and at very high prices. Continuous rebalancing
can create a strict local martingale from a true-martingale asset, and the
classification can reverse when rebalancing occurs on a finite grid. We
also establish limits on identification from price paths and examine what
liquidation rules, option contracts, and constituent information add to
the valuation problem.

The empirical work develops daily and intraday signals of return variance
rising with the price level. This pattern has a connection to bubble
mechanisms in both local-volatility and stochastic-volatility models.
The measurement problem is to distinguish it from changes in background
volatility. The daily method uses a predictable market volatility scale;
the intraday method compares price movements within each session while
updating its background state between sessions. We evaluate the methods
in simulations and historical records, with particular attention to how
calibration changes their alarm patterns.

The results show substantial sensitivity in the intended simulation
families and distinct behavior on market data. In shared-market-scale
and constant-elasticity simulations, the daily statistic detects 39--72\%
of the specified alternatives, compared with 7--14\% for a calibrated
explosive-price benchmark. It remains relatively quiet in crisis cohorts
and flags all four studied homebuilders about 400 trading sessions before
their 2005 peaks, versus one for the benchmark. The intraday score adds
out-of-sample drawdown information to that benchmark. These findings
support the signals as measurements of price--variance behavior, with
the theory specifying the additional assumptions needed to interpret
such behavior as a pricing bubble.

\section*{How to read the manuscript}
\phantomsection
\addcontentsline{toc}{section}{How to read the manuscript}
The four parts answer different questions within one framework.
Part~\ref{part:theory} defines the economic object and establishes what
rebalancing and observation can change. Part~\ref{part:signals} introduces the
proposed signals, their calibration, and their daily and intraday evidence.
Part~\ref{part:funds} follows the theory into finite trading and liquidation.
Part~\ref{part:options} asks what options and index constituents add to prices
alone. A reader interested primarily in monitoring can begin with
Section~\ref{sec:targets}, then read Part~\ref{part:signals}; the preceding
theory explains the limits on interpretation.

Section~\ref{sec:literature} introduces the relevant literature and explains
how the contributions build on it. The main text includes the short arguments that explain the results. Longer
derivations, implementation specifications, and the research-to-manuscript map
are in the appendices. Experiment numbers identify reproducible research
records, not a sequence the reader must follow. No new empirical experiment is
introduced in this manuscript: the tables and figures reorganize existing
results. The selected generated tables are read directly from the stored
experiment summaries at build time.

\clearpage
\tableofcontents
\clearpage
\pagenumbering{arabic}
\section{The questions and the contributions}\label{sec:targets}
An investor confronted with a rapidly appreciating stock might ask whether
the price is detached from fundamental value, whether volatility rises as
the price rises, whether a crash is becoming more likely, and whether a
leveraged position changes those risks. These questions concern different
aspects of the same episode. Understanding their connections helps determine
what to measure and how to interpret the result.

We focus on a particular pattern: return variance increasing with the price
level. It appears in models where the price is a strict local martingale,
and it offers something concrete to look for in market data. Measuring it
requires separating movements associated with the stock's price from changes
in market conditions and background volatility. The daily and intraday
signals developed here make that separation in different ways. Alongside
them, the theoretical analysis explains how the bubble concept changes
when we move from an asset to a rebalanced portfolio or an option contract.

Our starting point is the strict-local-martingale definition of a bubble
\cite{jarrow2007complete}. In discounted units, under a specified pricing
measure $Q$, a nonnegative local martingale $S$ is a supermartingale. Its
finite-horizon defect is
\begin{equation}\label{eq:defect}
 b_S(T)=S_0-\E_Q[S_T]\geq0.
\end{equation}
A positive defect means that $S$ fails to be a true martingale on that horizon.
The physical measure $P$ describes the distribution from which historical
observations are drawn. An expectation under $Q$ is a valuation object, not an
estimate of the investor's average realized return under $P$.

A widely used approach to monitoring bubbles looks for explosive price
growth: intervals on which prices behave like an autoregression with a
root greater than one. The generalized supremum augmented Dickey--Fuller
test (GSADF) is a prominent example \cite{phillips2012multiple}. We use an
empirically calibrated version as a benchmark. Its focus on price dynamics
provides a useful comparison with our focus on the relationship between
price and return variance. A third question, whether a signal predicts a future drawdown,
concerns outcomes under $P$. The literature already documents that large
run-ups and their characteristics can carry crash information without run-ups
necessarily predicting low average returns \cite{greenwood2017bubbles}.
Forecasting is therefore one way to examine usefulness, not the definition of
the object studied here.

\begin{table}[htbp]
\centering\small
\begin{tabularx}{\linewidth}{@{}L{0.22\linewidth}YY@{}}
\toprule
Object & Evidence it requires & What a positive finding means \\
\midrule
Pricing defect & A pricing law and a valuation convention & A shortfall of
terminal expectation relative to current value \\
Explosive behavior & A specified time-series null and calibration & Departure
toward explosive price dynamics in that testing framework \\
Price-level variance dependence & A variance model, nuisance treatment, and
calibration & An observed pattern better fitted by positive level dependence \\
Future loss risk & A forecast design and genuinely later outcomes & An
association with losses over a specified horizon \\
\bottomrule
\end{tabularx}
\caption{Four questions often grouped under the word ``bubble.'' A result in
one row does not automatically establish a result in another.}
\label{tab:objects}
\end{table}

\subsection{What this research adds}
The theoretical centerpiece is a classification of constant-proportion
wealth. An investor who continually restores a fixed fraction of wealth
to an asset can inherit its strictness or acquire strictness through a
different mechanism: rebalancing into an asset that can reach zero.
The result depends on exposure and on the asset's behavior at both ends
of its price range. Finite rebalancing and liquidation change the answer,
which gives the classification a direct interpretation for traded funds.
The result applies established martingale criteria; its prior-art comparison
is incomplete, and we make no separate novelty claim for the finite-trading
results.

The identification results explain what observed prices can contribute to
that theory. A finite record reveals behavior within a visited range, while
the martingale classification can depend on an unvisited tail. Stochastic
volatility adds another connection and another distinction. In Sin's model,
positive alignment of price and volatility shocks is central to strictness,
but different pricing measures can give the same observed law different
valuation implications. These results motivate measuring price--variance
association while keeping its statistical and valuation interpretations
explicit.

The methodological contribution is a pair of daily and intraday signals
built around that association. Both compare positive and nonpositive level
dependence using the same background-variance specification. The daily
method adjusts for recent market volatility, leaving the stock's own level
available to explain additional variation. The intraday method anchors
prices within each session and updates its background state between
sessions. Their construction addresses a recurring measurement problem:
a sufficiently flexible volatility adjustment can absorb the pattern we
want to detect.

The empirical contribution is an evaluation of those choices, including
simulation sensitivity, calibration, historical alarms, and a supplementary
forecasting comparison. In the shared-market-scale simulations and a
constant-elasticity diffusion family, the daily statistic detects roughly
39--72\% of the specified alternatives, against 7--14\% for calibrated
GSADF. In historical data, its sparse crisis-cohort alarms coexist with a
coherent early signal across D.R. Horton, Lennar, Pulte, and Toll Brothers,
all around 400 sessions before their 2005 peaks. The intraday score adds
modest out-of-sample drawdown information to GSADF. Part~\ref{part:signals}
presents these findings alongside the settings where each method is weaker.
Calibration is part of the result: the reference distribution determines
what counts as unusual and can materially change the resulting alarms.

\subsection{How to assess the evidence}
The principal distinction is between a pricing defect under $Q$ and a
measured pattern under $P$. Simulation detection rates have known generating
models; historical episode and calm-control labels are comparison groups.
Reported calibration and forecast results describe the evaluated procedures
and populations. Across those comparisons, GSADF remains the stronger general
run-up benchmark, the intraday score contributes incremental drawdown
information, and the daily statistic's contribution is its sensitivity and
selective historical alarm pattern rather than a demonstrated forecasting
gain.

The research used retrospective episodes and current-constituent universes,
with overlapping windows and repeated method development. Chronological
out-of-sample fits therefore test temporal transfer within this research
record, rather than an untouched research holdout. These shared qualifications
apply throughout the empirical part. More specific limitations are stated
where they change the interpretation of a result.
